\section*{Bab \ref{ch:b1:s-block} --- Hidrogen dan Blok s}

\begin{solution}{exo:b1:s-block:1}
\ce{KH} dan \ce{CaH2}: salin; \ce{SiH4} dan \ce{H2S}: kovalen;
\ce{PdH_{0.6}}: metalik. \ce{CaH2 + 2H2O -> Ca(OH)2 + 2H2}.
\end{solution}

\begin{solution}{exo:b1:s-block:2}
\ce{2Li + 2H2O -> 2LiOH + H2}, dan sama halnya dengan Na dan K;
\ce{Mg + 2HCl -> MgCl2 + H2}.
\end{solution}

\begin{solution}{exo:b1:s-block:3}
Basa: \ce{Na2O}, \ce{MgO} (\ce{MgO + 2HCl -> MgCl2 + H2O}). Asam:
\ce{SiO2}, \ce{SO3}, \ce{CO2} (\ce{SO3 + H2O -> H2SO4}). Amfoter:
\ce{Al2O3}, \ce{ZnO} (\ce{ZnO + 2OH- + H2O -> Zn(OH)4^2-}).
\end{solution}

\begin{solution}{exo:b1:s-block:4}
Pada pH 14 pasangan air adalah \ce{2H2O + 2e- -> H2 + 2OH-}, $E = -0.83$~V;
$\log K = 2(-0.83 + 2.71)/0.059 = 64$.
\end{solution}

\begin{solution}{exo:b1:s-block:5}
Di dalam nyala, tumbukan menaikkan sebuah elektron ke tingkat tereksitasi;
ketika kembali ke \omterm{def:b1:quantum-numbers:energy-level}{keadaan dasar}, atom itu memancarkan foton yang
energinya sama dengan selisih tingkat itu, yang khas untuk unsurnya. $E = hc/\lambda = 6.626 \times 10^{-34} \times
3.00 \times 10^8/(589 \times 10^{-9}) = \qty{3.375e-19}{J} = \qty{2.11}{eV}$.
% ledger: const:h, const:c, const:eV
\end{solution}

\begin{solution}{exo:b1:s-block:6}
$10^6/106.0 = \qty{9.43e3}{mol}$ abu soda: \qty{9.43e3}{mol} batu kapur,
\qty{0.944}{t}, dan $2 \times 9.43 \times 10^3/0.90 =
\qty{2.10e4}{mol}$ garam, \qty{1.23}{t}.
\end{solution}

\begin{solution}{exo:b1:s-block:7}
$[\ce{Mg^2+}] = 1.3/24.3 = \qty{0.053}{mol/L}$; $\mathrm{pH} = 14.00 -
\frac12(11.25 + \log 0.053) = 9.0$. \qty{53.5}{mol} magnesium dalam
\qty{1.0}{m^3} memerlukan \qty{53.5}{mol} \ce{Ca(OH)2}: \qty{4.0}{kg}.
\end{solution}

\begin{solution}{exo:b1:s-block:8}
Tetapan kesetimbangan (termodinamika) sangat besar untuk keduanya; laju
(kinetika) tidak ditentukan oleh $E^\circ$. Litium meleleh pada suhu yang
lebih tinggi, tetap padat dan tertutup lapisan produk yang lebih sukar
larut, dan melepaskan elektronnya kurang mudah daripada kalium (energi
ionisasinya lebih tinggi): litium bereaksi dengan tenang, kalium dengan
hebat.
\end{solution}

\begin{solution}{exo:b1:s-block:9}
$10^6/100.1 = \qty{9.99e3}{mol}$: CaO $9.99 \times 10^3 \times 56.1 =
\qty{0.560}{t}$; \ce{CO2} \qty{9.99e3}{mol}, $V = nRT/p = 9.99 \times 10^3
\times 8.314 \times 298.15/10^5 = \qty{248}{m^3}$.
\end{solution}

\begin{solution}{exo:b1:s-block:10}
$2.9 \times 10^8/6.9 = \qty{4.2e7}{kmol}$ Li (massa dalam kg), $\times 73.8/2 =
\qty{1.55e9}{kg}$: sekitar 1.6 juta ton litium karbonat.
$2.9 \times 10^8/8 = \num{3.6e7}$ baterai.
\end{solution}

\begin{solution}{exo:b1:s-block:11}
1.31, 0.84, 0.71 persen pada 20, 80, \qty{100}{\celsius}: \omterm{def:b1:precipitation:solubility}{kelarutannya}
menurun. Dalam keadaan panas, lebih sedikit karbonat yang tetap terlarut:
lebih banyak yang diperoleh kembali.
\end{solution}

\begin{solution}{exo:b1:s-block:12}
\ce{NaHCO3} adalah garam yang paling sukar larut di antara garam-garam
yang dapat dibentuk ion-ion itu, dan larutannya paling dahulu jenuh
terhadap garam itu. Amonia membuat larutan cukup basa untuk mengubah
\ce{CO2} menjadi \ce{HCO3-} sambil tetap dapat diperoleh kembali (sebagai
\ce{NH4Cl}, yang diregenerasi oleh kapur); natrium hidroksida akan habis
terpakai dan lebih mahal daripada produknya.
\end{solution}

\begin{solution}{pb:b1:s-block:1}
\textbf{1.} Li \qty{6.0}{g/L}, Mg \qty{2.4}{g/L}.
\textbf{2.} Empat meter kubik air garam menghasilkan satu: tiga meter
kubik air menguap.
\textbf{3.} Sinar matahari dan iklim yang kering memasok energinya
dengan cuma-cuma.
\textbf{4.} Natrium klorida, garam yang jauh paling melimpah, paling
dahulu jenuh.
\textbf{5.} Li $6000/6.9 = \qty{870}{mol}$; Mg $2400/24.3 = \qty{98.8}{mol}$.
\textbf{6.} Magnesium akan ikut mengendap bersama karbonat dan mencemari
produknya.
\textbf{7.} \ce{Mg^2+ + Ca(OH)2 -> Mg(OH)2 + Ca^2+}.
\textbf{8.} $98.8 \times 74.1 = \qty{7.32}{kg}$.
\textbf{9.} $[\ce{Mg^2+}] = \num{9.88e-5}$ mol/L yang tersisa: $\mathrm{pH} = 14.00 -
\frac12(11.25 - 4.01) = 10.38$.
\textbf{10.} Litium hidroksida mudah larut: tidak ada padatan yang
terbentuk pada pH ini.
\textbf{11.} \ce{Ca^2+ + CO3^2- -> CaCO3}.
\textbf{12.} $98.8 \times 106.0 = \qty{10.5}{kg}$.
\textbf{13.} \ce{2Li+ + CO3^2- -> Li2CO3}.
\textbf{14.} $435 \times 106.0 = \qty{46.1}{kg}$.
\textbf{15.} $435 \times 73.8 = \qty{32.1}{kg}$.
\textbf{16.} Litium karbonat lebih sukar larut dalam keadaan panas.
\textbf{17.} $0.0084 \times 1000 = \qty{8.4}{kg}$.
\textbf{18.} Dengan air panas dalam volume kecil, tempat zat itu paling
sukar larut.
\textbf{19.} $32.1 - 8.4 = \qty{23.7}{kg}$.
\textbf{20.} $23.7/32.1 = \qty{74}{\percent}$.
\textbf{21.} Cairan itu masih mengandung litium: cairan itu dikembalikan
ke kolam.
\textbf{22.} $2.9 \times 10^8/6.0 = \num{4.8e7}$ meter kubik.
\textbf{23.} \textbf{\qty{24}{kg}} (23.7) litium karbonat per meter
kubik.
\end{solution}
